% ---- pg-462 （书页 450） ----

文件有他们的签名. 但是另外一些文件没有签名. 因而人们不禁想知道: 这两人之中究竟是谁撰写了这些未署名文件? 很清楚, 我们没法去问他们. 但词的使用频数和句子的长度能帮助我们对此有争议的文件进行分类.

许多疾病只有通过非常昂贵的手术才能确诊. 通常, 我们希望通过某些容易观测但也可能误导的外部症候来诊断某种疾病. 这种方法有助于避免不必要的(而且是昂贵的)手术.

从以上例子中应该明白: 任何分类法则都不能提供准确无误的指派方法. 这是因为在所测量的总体特征之间可能不存在明确的界限. 也就是说, 两组特征可能是重叠的. 这样就可能将一个本属于 $\pi_2$ 的对象误分到 $\pi_1$, 或将属于 $\pi_1$ 的对象误分到 $\pi_2$.

\noindent\textbf{例 11.1}\quad (判别机动割草机的拥有者和非拥有者)

考虑某城市中的两群人: 拥有机动割草机的人群 $\pi_1$ 和不拥有机动割草机的人群 $\pi_2$. 在一次促销活动中, 某机动割草机的制造商为预测最佳销售前景, 想把该市家庭分为割草机的预期拥有者和非拥有者, 所依据的变量是: $x_1=$ 收入和 $x_2=$ 草坪大小. $n_1=12$ 个现有的拥有者与 $n_2=12$ 个现在的非拥有者两个随机样本所产生的数据列于\hyperref[tab:11.1]{表 11.1}.

\begin{mvtable}
\phantomsection\label{tab:11.1}\mvtabcap{表 11.1}
\begin{tabular}{@{}cccc@{}}
\toprule
\multicolumn{2}{c}{$\pi_1$: 机动割草机拥有者} & \multicolumn{2}{c}{$\pi_2$: 机动割草机非拥有者}\\
\cmidrule(r){1-2}\cmidrule(l){3-4}
$x_1$(收入 & $x_2$(草坪大小, & $x_1$(收入 & $x_2$(草坪大小,\\
单位: 千美元) & 单位: 千平方英尺) & 单位: 千美元) & 单位: 千平方英尺)\\
\midrule
90.0 & 18.4 & 105.0 & 19.6\\
115.5 & 16.8 & 82.8 & 20.8\\
94.8 & 21.6 & 94.8 & 17.2\\
91.5 & 20.8 & 73.2 & 20.4\\
117.0 & 23.6 & 114.0 & 17.6\\
140.1 & 19.2 & 79.2 & 17.6\\
138.0 & 17.6 & 89.4 & 16.0\\
112.8 & 22.4 & 96.0 & 18.4\\
99.0 & 20.0 & 77.4 & 16.4\\
123.0 & 20.8 & 63.0 & 18.8\\
81.0 & 22.0 & 81.0 & 14.0\\
111.0 & 20.0 & 93.0 & 14.8\\
\bottomrule
\end{tabular}
\end{mvtable}

这些数据示于\hyperref[fig:11.1]{图 11.1}. 我们看到, 割草机拥有者倾向于同时拥有较高的收入和较大的草坪, 虽然“收入”看上去是一个比“草坪大小”更好的判别量. 另一方面, 两个人群体有一定程度的重叠. 例如, 若将 $(x_1,x_2)$ 值落入区域 $R_1$(由图中的实线所界定)的人分配到 $\pi_1$, 而将 $(x_1,x_2)$ 值落入区域 $R_2$ 的人分配到 $\pi_2$, 将有可能犯某些错误. 某些拥有者将被错误地分类为非拥有者; 反之, 某些非拥有者也会被分类为拥有者. 基本的思路是建立一个法则(区域 $R_1$ 和 $R_2$), 使得犯上述错误的机会达到最小(见练习 11.2).\hfill$\blacksquare$

一个好的分类方法只应该产生少数错误分

\begin{mvfigure}
\mvinc{pg-462}{0.553,0.670,0.897,0.848}
\phantomsection\label{fig:11.1}\mvcap{图 11.1\quad 机动割草机的拥有者与非拥有者的收入和草坪大小}
\end{mvfigure}
